\section{Ограничения}
\label{sec:limitations}
\porozhdeno{\Claim{O01}{\istochnik{Если истинная таксономия неизвестна и положительного контроля нет, большее покрытие не позволяет заключить, что модель точнее~\cite{robeson2021}.}}
\Claim{O02}{Если обучение готовых моделей не контролировалось, наблюдаемое различие нельзя приписать исключительно референсной базе.}
\Claim{O03}{Если номенклатуры не согласованы, буквальное расстояние не измеряет чистое биологическое расхождение, а совпадение неназначенных профилей не подтверждает согласия известных таксонов.}
\Claim{O04}{Если утраченные при обработке риды Kr5 имеют иной состав, профиль сохранившихся \ProjectKrFiveRetentionPct{}\% не переносится на исходную пробу.}
\Claim{O05}{При неизвестных удалении праймеров Next\_515F/Next\_806R, регионе и условиях обучения моделей и номере запуска перенос результатов на другой протокол не подтверждён.}
\Claim{O06}{Если модельный отбор меняет множество ридов и знаменатель $N$, семейственные проценты нельзя трактовать как оценки на общем множестве, а при $N=0$ доля не определена.}
\Claim{O07}{Если образцы связаны общими точками отбора, химия измерена лишь в отдельных пробах КР1 и ЛЮБ1, а контаминация независимо не исключена, эти данные не обосновывают обобщения на районы Москвы или причинного объяснения профилей химическими свойствами почв.}}
